\section{General results: localisation, hierarchy and the finite-menu geometry}\label{sec:general}

\subsection{Two lemmas on quantifier order}

The first lemma is the abstract form of the separation; the second shows that it cannot be a static phenomenon.

\begin{lemma}[quantifier inclusion]\label{lem:quant}
Let $\{E_\lambda(z)\}_{\lambda\in\Lambda}$ be any family of plan sets. Then
\[
\Big\{z:\ \textstyle\bigcap_\lambda E_\lambda(z)\neq\varnothing\Big\}\ \subseteq\ \bigcap_\lambda\big\{z:\ E_\lambda(z)\neq\varnothing\big\},
\]
with equality exactly when, for every $z$ in the right-hand set, the family $\{E_\lambda(z)\}_\lambda$ admits a common selector. No topology on the plan space is needed; for a finite menu the check is a finite enumeration.
\end{lemma}

\begin{lemma}[the full cone is lossless at a fixed trajectory]\label{lem:cone}
For $v\in\mathbb R^n$: $v\ge0$ componentwise if and only if $w\cdot v\ge0$ for every $w\in\Wp$.
\end{lemma}

Proofs are in \ref{app:general}. Lemma~\ref{lem:cone} says that an assessor who checks the aggregate floor $w\cdot s\ge0$ at \emph{every} nonnegative weighting, along a \emph{fixed} trajectory, detects exactly the same violations as an assessor who checks each floor separately. Whatever separates the two doctrines must therefore enter through the quantifier over plans in \eqref{eq:central}, not through any blindness of scalarization at a profile.

\subsection{Hierarchy, localisation and weight-family monotonicity}

\begin{proposition}[hierarchy and localisation]\label{prop:hierarchy}
\textup{(i)} For every $w\in\Wp$ and every $z$,
\[
\Etyp(z)\ \subseteq\ E_w(z)\ \subseteq\ \Ephys(z)\ \subseteq\ \Eend(z).
\]
\textup{(ii)} For every $z$, $\Etyp(z)=\bigcap_{w\in\Wp}E_w(z)$; hence $\Vtyp=\{z:\ \bigcap_wE_w(z)\neq\varnothing\}$, and the central inclusion \eqref{eq:central} reads $\Vtyp\subseteq\Vweak$.
\end{proposition}

Part (i) is constraint-set monotonicity under nested safe sets, the standard monotonicity of viability kernels \citep{Aubin1991,Frankowska1989}. Part (ii) is the localisation that gives the paper its subject: the noncompensatory doctrine \emph{is} common-plan robustness over the full cone (Protocols 3 and 4 of Table~\ref{tab:protocols} coincide), so the only difference between the doctrines is the order of the two quantifiers.

\begin{proposition}[weight-family monotonicity]\label{prop:monotone}
For a weight family $W\subseteq\Wp$ let $\mathcal V_W=\bigcap_{w\in W}\mathcal V_w$. Then $\Vtyp\subseteq\mathcal V_W\subseteq\Vphys$, and $W_1\subseteq W_2$ implies $\mathcal V_{W_2}\subseteq\mathcal V_{W_1}$.
\end{proposition}

Restricting the weight family -- eliciting a narrower set of admissible weightings -- can only \emph{enlarge} the compensatory accepted set. The full cone is thus the strictest compensatory reading, and a separation exhibited at the full cone persists for every restricted family; the elicitation of weights, however careful, cannot remove it.

\subsection{The finite-menu geometric characterisation}

The mechanism behind the gap is geometric. Call a plan a \emph{margin-class} plan if its coordinate-wise admissibility is governed by a \emph{required-margin vector} $d_a\in\Rp$: plan $a$ is admissible at floor margins $s\ge0$ exactly when $d_a\le s$ componentwise, and it satisfies the destination constraint. A plan whose worst-case tube dips each floor $i$ by at most $d_{a,i}$ and returns is of this class; plans whose admissibility depends on another resource (reserve-financed plans) or that fail the destination are handled separately, as on the witness of Section~\ref{sec:witness}.

\begin{theorem}[finite-menu geometric characterisation]\label{thm:menu}
Let $\mathcal A$ be a finite menu of margin-class plans with required margins $\mathcal D=\{d_a:a\in\mathcal A\}$. Then
\begin{enumerate}[label=\textup{(\roman*)},leftmargin=2em]
\item coordinate-wise acceptance is $\{s:\ \exists a,\ d_a\le s\}=\bigcup_{a\in\mathcal A}(d_a+\Rp)$;
\item scalarized (weight-adaptive) acceptance is $\{s:\ \forall w\in\Wp\ \exists a,\ w\cdot d_a\le w\cdot s\}=\conv(\mathcal D)+\Rp$.
\end{enumerate}
Hence the acceptance gap is $\big[\conv(\mathcal D)+\Rp\big]\setminus\bigcup_{a}(d_a+\Rp)$ -- the part of the convex hull of the required margins that no single plan dominates -- and replacing the menu by its convex hull, i.e.\ admitting every convexified plan, collapses the gap to the empty set.
\end{theorem}

The proof (\ref{app:general}) is a separating-hyperplane argument; its content is displayed in Figure~\ref{fig:geometry}. Weight-adaptive acceptance of $s$ means that $s$ is not strictly separable from the up-closed convex hull of the required margins -- for every weighting some plan's margin is cheaper than $s$ in that weighting -- while coordinate-wise acceptance means that $s$ dominates some margin outright. The two coincide exactly when the up-closure of $\mathcal D$ is already convex, in particular whenever one plan dominates all others or the menu is convexified. Two consequences organise the rest of the paper: the gap is a property of the \emph{finiteness and nonconvexity of the menu}, not of either doctrine; and it is closed by convexifying the menu (Theorem~\ref{thm:blend}), whereas instruments that act on the weights alone cannot close it (Proposition~\ref{prop:monotone}). The nonconvexity at work is menu geometry; it is distinct from the Pareto-front nonconvexity under which weighted sums fail in multi-objective optimisation \citep{DasDennis1997}, which acts through frontier geometry under a single optimisation and has no quantifier over plans.

\begin{figure}[!htbp]
\centering
\includegraphics[width=\linewidth]{fig_finite_menu_geometry.png}
\caption{The finite-menu geometric characterisation (Theorem~\ref{thm:menu}) in the plane of two floor margins. Green: coordinate-wise acceptance, the union of the orthants $d_a+\mathbb R^2_+$ (some plan dominates). Black boundary: scalarized acceptance, $\conv(\mathcal D)+\mathbb R^2_+$. Hatched: the acceptance gap, where every weighting licenses some plan but no single plan keeps every floor. (a) A three-plan menu. (b) The margin-class plans of the witness datum, $d_{\textsc{fast}}=(2,0)$ and $d_{\textsc{slow}}=(0,2)$: the gap is the discrepancy triangle $\Q=\{s_1<2,\,s_2<2,\,s_1+s_2\ge2\}$ of Theorem~\ref{thm:witness}, with the witness point $(6/5,6/5)$ in its interior; the reserve-financed plan, which is outside the margin class, serves the triangle only when the reserve covers its cost.}
\label{fig:geometry}
\end{figure}

